\chapter{Biological Introduction}
\label{chap:bioIntr}
\section{From a Neural Tube to the Central Nervous System}
The whole central nervous system (\gls{CNS}), including the brain, brainstem and the spinal cord,
emerges from a single layered neural tube that is formed early in vertebrate development through the process of neurulation. 
By the time it closes (around embryonic day (\gls{E}) 9.5-10.5 in mice and E25-27 in humans), the neural tube already exhibits distinct anteroposterior and 
orsoventral patterning, giving rise to the three primary brain vesicles: the forebrain (Prosencephalon), midbrain (Mesencephalon), and hindbrain (Rhombencephalon) 
\citep{SAKAINeurulationMouseManner1989}. The forebrain is further divided into the Diencephalon and the telencephalon, which, once fully developed, comprises the cerebral cortex, including hippocampus and neocortex, as well as the basal ganglia and olfactory bulb.


%IMAGE BRAIN VESSEL
\begin{figure}[h]
    \centering
    \includegraphics[width=\textwidth]{brain_vessels.png}
    \caption{\textbf{Formation of the three primary brain vesicles during early neural development.} 
    The neural tube is already regionalized into distinct structures corresponding to the forebrain, midbrain, and hindbrain. \textcopyright Queensland Brain Institute, University of Queensland, Australia.}
    \label{fig:brainVessel}
\end{figure}
% -----------------------------------------------------

\section{Neurogenesis generates Neural Diversity}
With the closure of the neural tube, a series of tightly regulated processes initiate the formation of the complex brain. At this stage, the neural tube consists of just a single layer of neuroepithelial stem cells(\glspl{NSC}),
which must first undergo population expansion to generate sufficient progenitors for subsequent developmental stages. Between E8 and E10, 
NSCs undergo rapid symmetric division, a phase referred to as the expansion phase. Once an adequate number of NSCs has been established, neurogenesis begins. 
\\

Neurogenesis refers to the formation of neurons from NSCs and its neural progenitor cells (\glspl{NPC}) \citep{SUNGrowthFoldingMammalian2014}, and occurs throughout the whole neural tube. In rodents, neurogenesis spans approximately E10.5 to E18.5 \citep{GOTZCellBiologyNeurogenesis2005}, followed by radial neuronal migration, differentiation, as well as the formation of dendrites and axons for interconnection within the neural network \citep{KRIEGSTEINPatternsNeuronalMigration2004}.
\\

Neurogenesis is initiated by a switch from symmetric to asymmetric division of NECs. During these self-renewing asymmetric divisions, one daughter cell retains NEC identity, while the other differentiates into an apical radial glial cell, which we will refer to as apical progenitor (\gls{AP}) \citep{GOTZCellBiologyNeurogenesis2005}.


\begin{figure}[h]
    \centering
    \includegraphics[width=\textwidth]{neurogenesis.png}
    \caption{\textbf{Development of the mouse neocortex.} Schematic overview of neocortical development in the mouse. During early embryogenesis (E9.5–E11.5), 
    neuroepithelial cells (NECs) in the ventricular zone (\gls{VZ}) 
    undergo symmetric proliferative divisions, expanding the progenitor 
    pool. Around E11.5, NECs transition into apical progenitors (APs), 
    which initially self-renew and subsequently enter a neurogenic phase. 
    Between E11.5 and E13.5, APs generate early-born deep-layer neurons 
    (DLNs), which migrate basally toward the cortical plate. 
    From approximately E13.5 onward, APs also produce basal progenitors 
    (BPs) that populate the subventricular zone (SVZ) and give rise to late-born upper-layer neurons (ULNs). Blood vessels invade radially from the basal lamina toward the apical surface, suggesting the establishment of an oxygen gradient decreasing from basal to apical regions. Figure created with BioRender.com.}
    \label{fig:neurogenesis}
\end{figure}


\subsection{Apical Progenitors}
Both, NECs and APs reside within and divide in the same germinal zone, the ventricular zone (\gls{VZ}), which is defined by its direct contact with ventricles that are cerebrospinal fluid filled on the apical side. APs exhibit several characteristic features: (1) an elongated 
morphology with both apical and basal processes, allowing direct contact with the ventricular and pial surfaces. (2) interkinetic nuclear migration (INM), a process in which the nucleus moves from the apical toward the basal side during the cell cycle, while mitosis occurs at the ventricular surface and (3) the expression of astroglial marker genes such as nuclear Pax6 (Paired box 6), Sox2 (sex determining region Y-box 2) \citep{GOTZPax6ControlsRadial1998}, glutamate transporter (GLAST) (Shibata et al. 1997), glial fibrillary acidic protein (GFAP)
\citep{LEVITTImmunoperoxidaseLocalizationGlial1980}, and  glutamine synthetase (\gls{GS}) \citep{ENGLUNDPax6Tbr2Tbr12005}.
\\

APs represent a distinct and versatile progenitor population, capable of adopting multiple cell fates during neocortex development. Following an initial expansion phase ($\sim$E10-E12) , APs undergo a progressive shift toward asymmetric cell division. In the first phase of asymmetry divisions ($\sim$E12-E13), APs generate postmitotic neurons destined for deep cortical layers (deep-layer neuron, \gls{DLN}). \newline
\\

As neurogenesis proceeds ($\sim$E13-E16), their asymmetric division increasingly gives rise to more differentiated progenitor populations: basal progenitors (\glspl{BP}).
In comparison to NECs and APs, BPs detach from the apical surface and migrate outward to establish the secondary germinal layer, called the subventricular zone (\gls{SVZ}).


\subsection{Basal Progenitors}

Basal progenitors represent a heterogeneous secondary progenitor class that are more difficult to clearly separate from other cell types. While some BPs fully detach from the ventricular surface, others retain a basal process and maintain a gene expression profile more closely related to their AP origin \citep{pilzAmplificationProgenitorsMammalian2013}. 
Basal radial glia, for example, represents a minority of BPs in rodents, accounting for approximately 5\% \citep{MARTINEZ-CERDENOComparativeAnalysisSubventricular2012}, and are still attached to the basal surface. 
\\

Basal intermediate progenitors on the other hand, represent a major subclass in BPs and are no longer in contact with either the apical or basal surface. Furthermore, basal intermediate progenitors exhibit limited proliferative capacity in comparison to APs. They typically undergo a single proliferative division, before producing two upper layer neurons (\glspl{ULN}) \citep{THORIndirectNeurogenesisSpace2024}. 
In addition, they are characterized by different marker gene expression, including Tbr2/Eomes (T-box brain protein 2/Eomesodermin) \citep{ENGLUNDPax6Tbr2Tbr12005}. Here, BPs refers specifically to the subclass of basal intermediate progenitors. 

\subsection{Migration Of Upper-Layer And Deep-Layer Neurons}

DLNs and ULNs originate from distinct progenitor populations within the developing neocortex. DLNs are primarily generated by APs in the VZ, while ULNs arise from BPs in the SVZ. Regardless of their origin, both neuronal populations must undergo radial migration towards the pial surface in order to reach their final positions in the deeper and upper layers of the neocortex, respectively \citep{NADARAJAHModesNeuronalMigration2002}.
\\

Two principal modes of radial migration have been described. In early stages, neurons predominantly undergo somal translocation, a process in which the neurons extend a basal process that anchors to the pial surface while the nucleus moves along this process toward the basal side, resulting in progressive shortening of the process while maintaining attachment. At later stages, neurons primarily adopt a locomotion-based mode of migration, also referred to as gliophilic migration, in which they migrate along radial glial fibers of APs \citep{NADARAJAHModesNeuronalMigration2002}.
Through radial migration along the fibers of APs, DLNs settle near the pial surface and form the initial cortical layer. As neurogenesis progresses, ULNs migrate past these earlier-born cells to establish progressively more superficial layers \citep{BUCHSBAUMNeuronalMigrationCNS2019}.
\\

This sequential pattern of migration and positioning results in the characteristic inside-out organization of the cortex, whereby early-born neurons occupy deeper layers and later-born neurons populate the upper layers.

